\documentclass[11pt,a4paper]{article}
% Estilo comum dos PDFs individuais da Sprint 1.
\usepackage[utf8]{inputenc}
\usepackage[T1]{fontenc}
\usepackage[brazilian]{babel}
\usepackage[a4paper,margin=2.35cm]{geometry}
\usepackage{xcolor}
\usepackage{hyperref}
\usepackage{enumitem}
\usepackage{listings}
\usepackage{parskip}
\usepackage{needspace}
\definecolor{azul}{RGB}{20,63,112}
\hypersetup{colorlinks=true,linkcolor=azul,urlcolor=azul}
\setlist[itemize]{leftmargin=*,itemsep=3pt}
\lstset{basicstyle=\ttfamily\small,frame=single,breaklines=true,columns=fullflexible}
\sloppy
\newcommand{\codigo}[1]{\texttt{#1}}
\newcommand{\titulo}[2]{\begin{center}{\LARGE\bfseries #1\par}\vspace{4pt}{\large #2\par}\vspace{8pt}\rule{0.8\textwidth}{0.4pt}\par\vspace{5pt}{\small Matheus Moreli Durão Silva -- Matrícula 12621CBS217\par Cibersegurança - Bacharelado - Noturno\par}\end{center}}
\newcommand{\cardfonte}[1]{\par\smallskip\noindent\textit{Referência: }\href{#1}{Trello - enunciado e checklist do card}.}

\begin{document}\titulo{S1\_12 - Dicionários}{Vincular vulnerabilidades a um ativo}
\section*{Entregável correto}
O exercício solicitado registra e atualiza vulnerabilidades associadas a ativos usando dicionários. Isso já se aproxima do modelo do sistema final: em vez de uma lista de textos soltos, cada ativo possui campos com nomes.
\section*{Estrutura e operações}
Um dicionário relaciona \textbf{chave} e \textbf{valor}: \codigo{ativo = \{"nome": "Servidor", "vulnerabilidades": []\}}. O acesso é feito pela chave, como \codigo{ativo["nome"]}. Para inserir ou atualizar, atribui-se: \codigo{ativo["criticidade"] = "Alta"}. Para remover, usa-se \codigo{del ativo["criticidade"]} ou \codigo{pop}. \codigo{keys()}, \codigo{values()} e \codigo{items()} permitem percorrer dados.
\begin{lstlisting}
ativo = {"nome": "Servidor", "vulnerabilidades": []}
ativo["vulnerabilidades"].append("porta 22 exposta")
ativo["vulnerabilidades"][0] = "porta SSH exposta"
for chave, valor in ativo.items():
    print(f"{chave}: {valor}")
\end{lstlisting}
\section*{Por que há uma lista dentro do dicionário?}
Um ativo tem uma só chave \codigo{nome}, mas pode ter muitas vulnerabilidades. Por isso, o valor de \codigo{vulnerabilidades} é uma lista. Essa combinação de estruturas modela o mundo real: dicionário descreve um registro e lista guarda ocorrências repetidas ligadas a ele.
\section*{Aplicação no projeto}
O \codigo{app.py} usa dicionários para representar ativos, e a vulnerabilidade fica ligada ao ativo correto. Atualizar um valor por chave é a base do menu de edição. Ao salvar em JSON, a mesma forma de dados continua legível.
\section*{Checklist demonstrado}
Chave/valor, criação/acesso, inserção/atualização/remoção, \codigo{keys/values/items}, laços e modelagem foram demonstrados. Teste: cadastrar uma vulnerabilidade, atualizar a descrição e listar o dicionário para confirmar a alteração.
\cardfonte{https://trello.com/c/N1J9NE4k}
\end{document}
